\sectionheading{5.20 সমস্যাটি বিশ্লেষণ ও IPO মডেল (Problem Analysis and IPO Model)}

কোনো Program লেখার আগে সবচেয়ে গুরুত্বপূর্ণ কাজ হলো
\textbf{সমস্যাটি সঠিকভাবে বোঝা}।

Problem কী চাচ্ছে, কী Data দেওয়া হবে, কোন Processing করতে হবে
এবং কী Result পাওয়া যাবে—এসব নির্ধারণ করাই
Problem Analysis-এর মূল কাজ।

%================================================
% Problem Analysis Definition
%================================================
\booksubheading{Problem Analysis কাকে বলে?}

\begin{definitionbox}

\textbf{\color{darkTeal}সংজ্ঞা}

কোনো Problem-এর উদ্দেশ্য, প্রয়োজনীয় Input, করণীয় Processing,
প্রত্যাশিত Output এবং প্রয়োজনীয় শর্তসমূহ চিহ্নিত ও বিশ্লেষণ
করার প্রক্রিয়াকে \textbf{Problem Analysis} বলে।

\end{definitionbox}

\textbf{সহজভাবে মনে রাখবে—}

\begin{center}

\textbf{
Problem কী?
\(\longrightarrow\)
কী দেওয়া আছে?
\(\longrightarrow\)
কী করতে হবে?
\(\longrightarrow\)
কী ফল চাই?
}

\end{center}

%================================================
% Main Tasks
%================================================
\booksubheading{Problem Analysis-এর প্রধান কাজ}

একটি Programming Problem বিশ্লেষণের সময় সাধারণত—

\begin{compacttwocolpoints}

\item Problem-এর \textbf{মূল উদ্দেশ্য} নির্ধারণ করা হয়।

\item কী কী \textbf{Input} প্রয়োজন তা চিহ্নিত করা হয়।

\item Input-এর ওপর কী \textbf{Process বা Calculation} হবে তা ঠিক করা হয়।

\item কী \textbf{Output} দিতে হবে তা নির্ধারণ করা হয়।

\item কোনো \textbf{Condition বা Constraint} আছে কি না তা দেখা হয়।

\item Solution-এর কাজগুলো কোন ক্রমে সম্পন্ন হবে তা বিবেচনা করা হয়।

\end{compacttwocolpoints}

এগুলো পরিষ্কার হলে Algorithm, Flowchart এবং Program Code
তৈরি করা সহজ হয়।

%================================================
% Importance
%================================================
\booksubheading{Problem Analysis কেন প্রয়োজন?}

সঠিকভাবে Problem Analysis না করলে Program-এর Syntax ঠিক হলেও
Solution ভুল হতে পারে।

Problem Analysis-এর মাধ্যমে—

\begin{compacttwocolpoints}

\item Problem সম্পর্কে পরিষ্কার ধারণা পাওয়া যায়।

\item প্রয়োজনীয় ও অপ্রয়োজনীয় Data আলাদা করা যায়।

\item Input ও Output সহজে নির্ধারণ করা যায়।

\item প্রয়োজনীয় Formula বা Condition চিহ্নিত করা যায়।

\item Algorithm ও Flowchart তৈরি সহজ হয়।

\item ভুল Solution তৈরির সম্ভাবনা কমে।

\end{compacttwocolpoints}

%================================================
% IPO Model
%================================================
\booksubheading{IPO Model}

Program-এর Problem বিশ্লেষণের একটি সহজ পদ্ধতি হলো
\textbf{IPO Model}।

\begin{center}
\begin{adjustbox}{max width=0.98\linewidth}
\fcolorbox{darkBlue}{softBlue}{
    \parbox[c][9mm][c]{0.18\linewidth}{
        \centering\textbf{Input}
    }
}
\quad
\(\longrightarrow\)
\quad
\fcolorbox{mainTeal}{softGreen}{
    \parbox[c][9mm][c]{0.20\linewidth}{
        \centering\textbf{Process}
    }
}
\quad
\(\longrightarrow\)
\quad
\fcolorbox{darkBlue}{softBlue}{
    \parbox[c][9mm][c]{0.18\linewidth}{
        \centering\textbf{Output}
    }
}
\end{adjustbox}
\end{center}

\begin{center}

\textbf{
IPO = Input -- Process -- Output
}

\end{center}

আগের Program Organization Topic-এ Input, Process ও Output-এর
মৌলিক ধারণা আলোচনা করা হয়েছে।

এখানে এগুলোকে \textbf{Problem Analysis-এর একটি Model}
হিসেবে ব্যবহার করা হবে।

%================================================
% IPO Definition
%================================================
\booksubheading{IPO Model কাকে বলে?}

\begin{definitionbox}

\textbf{\color{darkTeal}সংজ্ঞা}

কোনো Problem-এর Solution নির্ধারণের জন্য প্রয়োজনীয়
\textbf{Input}, সেই Input-এর ওপর সম্পাদিত \textbf{Process}
এবং প্রত্যাশিত \textbf{Output} চিহ্নিত করে Problem-কে
বিশ্লেষণের পদ্ধতিকে \textbf{IPO Model} বলা হয়।

\end{definitionbox}

%================================================
% IPO Components
%================================================
\booksubheading{IPO Model-এর তিনটি অংশ}

\begin{tabularx}{\linewidth}{|>{\bfseries}p{0.22\linewidth}|Y|}
\hline

\rowcolor{softBlue}
\multicolumn{1}{|c|}{\textbf{অংশ}}
&
\multicolumn{1}{c|}{\textbf{Problem Analysis-এ প্রশ্ন}}
\\
\hline

Input
&
কী কী Data দেওয়া হবে বা গ্রহণ করতে হবে?
\\
\hline

Process
&
Data-এর ওপর কী কাজ, Formula বা Calculation হবে?
\\
\hline

Output
&
শেষে কী Result পাওয়া বা প্রদর্শন করা হবে?
\\
\hline

\end{tabularx}

\begin{center}

\textbf{
কী নেব?
\(\longrightarrow\)
কী করব?
\(\longrightarrow\)
কী দেখাব?
}

\end{center}

%================================================
% Example 1
%================================================
\begin{examplebox}{উদাহরণ ১: দুটি সংখ্যার যোগফল}

Problem: দুটি সংখ্যার যোগফল নির্ণয় করতে হবে।

\begin{tabularx}{\linewidth}{|>{\bfseries}p{0.22\linewidth}|Y|}
\hline

\rowcolor{softBlue}
\multicolumn{1}{|c|}{\textbf{IPO অংশ}}
&
\multicolumn{1}{c|}{\textbf{বিশ্লেষণ}}
\\
\hline

Input
&
দুটি সংখ্যা \textenglish{\ttfamily a, b}
\\
\hline

Process
&
\textenglish{\ttfamily sum = a + b}
\\
\hline

Output
&
\textenglish{\ttfamily sum}
\\
\hline

\end{tabularx}

\vspace{1mm}

\begin{center}

\textbf{
a, b
\(\longrightarrow\)
a + b
\(\longrightarrow\)
Sum
}

\end{center}

\end{examplebox}

%================================================
% Example 2
%================================================
\begin{examplebox}{উদাহরণ ২: আয়তক্ষেত্রের ক্ষেত্রফল}

Problem: দৈর্ঘ্য ও প্রস্থ দেওয়া থাকলে আয়তক্ষেত্রের ক্ষেত্রফল
নির্ণয় করতে হবে।

\begin{tabularx}{\linewidth}{|>{\bfseries}p{0.22\linewidth}|Y|}
\hline

\rowcolor{softBlue}
\multicolumn{1}{|c|}{\textbf{IPO অংশ}}
&
\multicolumn{1}{c|}{\textbf{বিশ্লেষণ}}
\\
\hline

Input
&
দৈর্ঘ্য \textenglish{\ttfamily l}, প্রস্থ \textenglish{\ttfamily w}
\\
\hline

Process
&
\textenglish{\ttfamily area = l * w}
\\
\hline

Output
&
ক্ষেত্রফল \textenglish{\ttfamily area}
\\
\hline

\end{tabularx}

এখানে Problem Analysis করার সময় প্রয়োজনীয় Formula-টিও
চিহ্নিত করা হয়েছে।

\end{examplebox}

%================================================
% Example 3
%================================================
\begin{examplebox}{উদাহরণ ৩: সংখ্যা জোড় না বিজোড়}

Problem: একটি সংখ্যা Even না Odd নির্ণয় করতে হবে।

\begin{tabularx}{\linewidth}{|>{\bfseries}p{0.22\linewidth}|Y|}
\hline

\rowcolor{softBlue}
\multicolumn{1}{|c|}{\textbf{IPO অংশ}}
&
\multicolumn{1}{c|}{\textbf{বিশ্লেষণ}}
\\
\hline

Input
&
একটি পূর্ণসংখ্যা \textenglish{\ttfamily n}
\\
\hline

Process
&
\textenglish{\ttfamily n \% 2 == 0} পরীক্ষা
\\
\hline

Output
&
Even অথবা Odd
\\
\hline

\end{tabularx}

এখানে শুধু Calculation নয়, একটি
\textbf{Condition}-ও Problem Analysis-এর অংশ।

\end{examplebox}

%================================================
% Relationship
%================================================
\booksubheading{IPO Model ও Problem Analysis-এর সম্পর্ক}

Problem Analysis একটি বিস্তৃত প্রক্রিয়া।
IPO Model সেই বিশ্লেষণকে সহজভাবে তিনটি প্রধান ভাগে সাজাতে সাহায্য করে।

\begin{center}

\textbf{
Problem Analysis
\(\longrightarrow\)
Input শনাক্ত
\(\longrightarrow\)
Process নির্ধারণ
\(\longrightarrow\)
Output নির্ধারণ
}

\end{center}

\begin{tabularx}{\linewidth}{|>{\bfseries}p{0.28\linewidth}|Y|}
\hline

\rowcolor{softBlue}
\multicolumn{1}{|c|}{\textbf{বিষয়}}
&
\multicolumn{1}{c|}{\textbf{মূল ধারণা}}
\\
\hline

Problem Analysis
&
পুরো Problem বুঝে উদ্দেশ্য, Data, Condition ও Solution নির্ধারণ করা।
\\
\hline

IPO Model
&
Problem-এর Input, Process ও Output সুশৃঙ্খলভাবে সাজানো।
\\
\hline

\end{tabularx}

%================================================
% Computational Thinking Relationship
%================================================
\booksubheading{Computational Thinking-এর সঙ্গে সম্পর্ক}

আগের Topic-এর Computational Thinking-এর সঙ্গে
Problem Analysis-এর ঘনিষ্ঠ সম্পর্ক রয়েছে।

\begin{tabularx}{\linewidth}{|>{\bfseries}p{0.29\linewidth}|Y|}
\hline

\rowcolor{softBlue}
\multicolumn{1}{|c|}{\textbf{ধারণা}}
&
\multicolumn{1}{c|}{\textbf{Problem Analysis-এ ভূমিকা}}
\\
\hline

Decomposition
&
Problem-কে ছোট অংশে ভাগ করতে সাহায্য করে।
\\
\hline

Abstraction
&
প্রয়োজনীয় Data নির্বাচন ও অপ্রাসঙ্গিক তথ্য বাদ দিতে সাহায্য করে।
\\
\hline

Algorithmic Thinking
&
Solution-এর ধাপ নির্ধারণে সাহায্য করে।
\\
\hline

IPO Model
&
Input, Process এবং Output পরিষ্কারভাবে নির্ধারণ করে।
\\
\hline

\end{tabularx}

এগুলো একই বিষয় নয়; বরং Problem Solving-এর বিভিন্ন সহায়ক ধারণা।

%================================================
% Checklist
%================================================
\booksubheading{Problem Analysis-এর সহজ Checklist}

কোনো নতুন Programming Problem দেখলে নিচের পাঁচটি প্রশ্ন করো—

\begin{enumerate}

\item Problem আমাকে \textbf{কী বের করতে বলেছে?}

\item কী কী \textbf{Input} দেওয়া আছে বা নিতে হবে?

\item কোনো \textbf{Formula বা Condition} আছে কি?

\item Input-এর ওপর কী \textbf{Process} করতে হবে?

\item শেষ পর্যন্ত কী \textbf{Output} দেখাতে হবে?

\end{enumerate}

এই প্রশ্নগুলোর উত্তর পরিষ্কার থাকলে পরবর্তী Algorithm,
Flowchart এবং Program তৈরি অনেক সহজ হয়।

%================================================
% Final Reminder
%================================================
\begin{rememberbox}

\textbf{\color{deepOrange}মনে রাখবে}

\begin{itemize}

\item Program লেখার আগে \textbf{Problem Analysis} করা উচিত।

\item Problem Analysis-এ উদ্দেশ্য, Input, Process, Output এবং
প্রয়োজনীয় Condition চিহ্নিত করা হয়।

\item \textbf{IPO = Input -- Process -- Output}।

\item IPO Model হলো Problem Analysis-এর একটি সহজ পদ্ধতি।

\item \textbf{Input} = কী Data নেওয়া হবে।

\item \textbf{Process} = Data-এর ওপর কী কাজ হবে।

\item \textbf{Output} = কী Result পাওয়া যাবে।

\item Formula ও Condition-ও Problem Analysis-এর সময়
নির্ধারণ করা হতে পারে।

\item দ্রুত মনে রাখার ধারা—

\begin{center}

\textbf{
Problem বুঝি
\(\longrightarrow\)
Input খুঁজি
\(\longrightarrow\)
Process ঠিক করি
\(\longrightarrow\)
Output নির্ধারণ করি
}

\end{center}

\end{itemize}

\end{rememberbox}

\small

%================================================
% 5.21 Program Development
%================================================
